\section*{Index restrictions and switching (author source pp.22:12--22:14)}
We create a distribution on restrictions on the variables of $R^\ell\Ref(F,s)$ as follows: For every index $u$ and every $i\in[\ell]$, we set $P_i[u]$ to0 or1, with probability $1/2$ for each value. Let $U$ be the set of indices such that $P_i[u]=1$ for all $i\in[\ell]$. For each variable $x$ of $R^\ell\Ref(F,s)$ not of the form $P_i[u]$ mentioning an index in $U$, we set $x$ to0 or1, with probability $1/2$ for each value.
For a decision tree $T$ querying variables of $\Ref(F,s)$, we define the index-height of $T$ as the maximum number of indices mentioned by variables over all branches that do not falsify axioms of $\Ref(F,s)$. For a formula $G$, We denote by $\hbar(G)$ the minimum index-height of a decision tree representing $G$.
The following lemma is from[10]. We give a proof because in[10] the lemma is stated not for $R^\ell\Ref(F,s)$ but a variation, plus we view the following proof to be simpler.
\begin{originalresult}{Lemma 5.4 [10]}
Let $F$ be a CNF formula in $n$ variables,$k$ and $\ell$ integers with $0<k\le\ell$, and $G$ a $k$-DNF formula over the variables of $R^\ell\Ref(F,s)$, where $s\le2^{\delta n}$ for some $\delta<1$. Then for large enough $n$,
\[\mathbb P[\hbar(G|_\rho)>h]\le2^{-\frac h{n^{k-1}}\gamma(k)},\]
where $\gamma(0)=1,\gamma(i)=(\log e)(i4^{i+1})^{-1}\gamma(i-1)$.
\end{originalresult}
\begin{proof}
Let $h_i:=h\gamma(i-1)/(4n^{i-1})$. We will show, by induction on $k$, that for every $k$ and $\ell$ with $k\le\ell$, for every $k$-DNF formula $G$ over the variables of $R^\ell\Ref(F,s)$,
\[\mathbb P\left[\hbar(G|_\rho)>\sum_{i=1}^kh_i\right]\le2^{-\frac h{n^{k-1}}\gamma(k)}\]
for large enough $n$.
If $k=0$, $F$ is a constant and can be represented by a decision tree of height0. Suppose $k>0$. We call $G$ wide if there are at least $h_k/k$ terms in $G$ over disjoint sets of indices, and call $G$ narrow otherwise. Suppose $G$ is wide. A literal in a term $t$ of $G$ is satisfied with probability at least $1/4$: Literals on a variable $P_i[u]$ are satisfied with probability $1/2$. For any other literal $x^\epsilon$ of $t$ mentioning the index $u$, since $k\le\ell$, there must be a variable $P_i[u]$ not in $t$, which is made0 with probability $1/2$, in which case $x^\epsilon$ will be satisfied with probability $1/2$. Hence
\begin{align*}
\mathbb P[\hbar(G|_\rho)>h]&\le\mathbb P[G|_\rho\ne1]\le(1-4^{-k})^{\frac{h\gamma(k-1)}{4kn^{k-1}}}\\
&\le2^{-\frac h{n^{k-1}}(\log e)(k4^{k+1})^{-1}\gamma(k-1)}=2^{-\frac h{n^{k-1}}\gamma(k)}.
\end{align*}
Suppose now that $G$ is narrow. Take a maximal set of terms over disjoint sets of indices, and let $H$ be the set of indices that are mentioned by the terms of this set. Notice that $|H|\le h_k$ and that every term of $G$ contains some variable (or its negation) that mentions an index in $H$. Let $\pi$ be a restriction that
\begin{enumerate}\item sets all variables mentioning an index in $H$ and leaves all other variables unset, and\item does not falsify any axioms of $R^\ell\Ref(F,s)$.\end{enumerate}
The second condition means in particular that if $U$ is the set of indices $u$ for which $\pi$ sets $P_i[u]$ to1 for all $i$, then for all $u\in U$, there will be exactly one $v$ such that $L[u,v]$ is true, exactly one $v$ such that $R[u,v]$ is true, exactly one $i$ such that $V[u,i]$ is true, and exactly one $j$ such that $I[u,j]$ is true, making the total number of such $\pi$'s to be at most
\[S^{|U|}2^{(|H|-|U|)n_0}\]
where $S:=s^2(n+1)(m+1)2^{2n}$ and $n_0$ is the number of variables of $R^\ell\Ref(F,s)$ mentioning a fixed index $u$.
Let $A_\pi$ be the event that $\pi$ is consistent with $\rho$ and $\hbar((G|_\rho)|_\pi)>\sum_{i=i}^{k-1}h_i$. We have that
\begin{align*}
\mathbb P[A_\pi]&=\mathbb P\left[\hbar((G|_\rho)|_\pi)>\sum_{i=i}^{k-1}h_i\mid\rho\text{ con.with }\pi\right]\mathbb P[\rho\text{ con.with }\pi]\\
&=\mathbb P\left[\hbar((G|_\pi)|_\rho)>\sum_{i=i}^{k-1}h_i\right]\mathbb P[\rho\text{ con.with }\pi]\\
&\le2^{-\frac h{n^{k-2}}\gamma(k-1)}2^{-\ell|U|}2^{-(|H|-|U|)n_0}.
\end{align*}
Hence, we get
\begin{align*}
\mathbb P\left[\bigcup_\pi A_\pi\right]&\le\sum_\pi\mathbb P[A_\pi]\\
&\le\sum_{U\subseteq H}S^{|U|}2^{(|H|-|U|)n_0}2^{-\frac h{n^{k-2}}\gamma(k-1)}2^{-\ell|U|}2^{-(|H|-|U|)n_0}\\
&=\sum_{r=0}^{|H|}\binom{|H|}{r}S^r2^{-\frac h{n^{k-2}}\gamma(k-1)}2^{-\ell r}\\
&=(S/2^\ell+1)^{|H|}2^{-\frac h{n^{k-2}}\gamma(k-1)}\\
&\le S^{|H|}2^{-\frac h{n^{k-2}}\gamma(k-1)}.
\end{align*}
Since $s\le2^{\delta n}$ for some $\delta<1$, the quantity
\[S^{|H|}=\left(s^2(n+1)(m+1)2^{2n}\right)^{\frac h{4n^{k-1}}\gamma(k-1)}\]
will be at most $2^{\frac{\varepsilon h}{n^{k-2}}\gamma(k-1)}$ for some $\varepsilon<1$ for large enough $n$, therefore
\[\mathbb P\left[\bigcup_\pi A_\pi\right]\le2^{-\frac h{n^{k-1}}\gamma(k)}\]
for large enough $n$.
In the event $\left(\bigcup_\pi A_\pi\right)^c$, that is the event that for every $\pi$ consistent with $\rho$, $\hbar((G|_\rho)|_\pi)\le\sum_{i=1}^{k-1}h_i$, we can construct a decision tree for $G|_\rho$ of index-height at most $\sum_{i=1}^kh_i$ as follows: We first query all variables mentioning an index in $H$ left unset by $\rho$. Then, at each branch $\pi$ of the resulting tree, we append a decision tree of minimum index-height representing $(G|_\rho)|_\pi$.
\end{proof}
